\section{Discussão}

\textbf{Convergência versus estabilidade.} A taxa de aprendizado $\alpha$ mostrou-se o hiperparâmetro mais determinante. Com $\alpha = 0{,}001$, a perda estagnou em $2{,}3016$ --- próxima à entropia uniforme $\ln(10) \approx 2{,}3026$ --- com acurácia de $12{,}71\%$ (nível do acaso); já $\alpha = 0{,}01$ atingiu apenas $30{,}61\%$ (perda $2{,}2840$). Como o \textit{full-batch} executa uma única atualização por época e os pesos iniciais são pequenos ($\sigma = 0{,}01$), os gradientes iniciais têm baixa magnitude e $200$ passos são insuficientes. Elevar $\alpha$ para $0{,}1$ e $0{,}5$ acelerou fortemente a convergência, atingindo $86{,}11\%$ (perda $0{,}5320$) e $92{,}43\%$ (perda $0{,}2585$, em $280{,}6$~s), embora taxas excessivas induzissem oscilações ou divergência.

\textbf{Retornos decrescentes e orçamento computacional.} Variar o número de épocas (com $\alpha = 0{,}1$ e $h = 64$) evidencia ganhos marginais decrescentes. Dobrar de $50$ para $100$ épocas reduziu a perda de $2{,}1144$ para $1{,}0912$ e elevou a acurácia em $+32{,}00$ pontos percentuais (p.p.), de $44{,}28\%$ ($113{,}9$~s) para $76{,}28\%$ ($196{,}4$~s). De $100$ para $200$ épocas, o ganho caiu para $+9{,}83$~p.p. ($86{,}11\%$, $355{,}5$~s). Adicionar mais $300$ épocas ($200 \to 500$) rendeu apenas $+4{,}10$~p.p. ($90{,}21\%$, perda $0{,}3417$), mais que dobrando o tempo ($355{,}5$~s $\to$ $729{,}3$~s).

\textbf{Capacidade da rede versus custo e generalização.} Com $h = 4$ ($3.190$ parâmetros), ocorre claro \textit{underfitting} ($69{,}26\%$, perda $0{,}9421$, $84{,}7$~s) pelo gargalo informacional para $10$ classes. Como os parâmetros escalam linearmente ($784h + h + 10h + 10$), expandir $h$ para $16$, $64$ e $256$ eleva o total para $12.730$, $50.890$ e $203.530$ parâmetros e o tempo de $115{,}9$~s para $355{,}5$~s e $922{,}4$~s. Apesar do avanço até $h = 64$ ($83{,}74\% \to 86{,}11\%$), quadruplicar para $h = 256$ adicionou apenas $+1{,}52$~p.p. ($87{,}63\%$, perda $0{,}4755$), elevando o risco de \textit{overfitting} sob treinos mais longos.

\textbf{Limitações arquiteturais e algorítmicas.} Três fatores restringem o modelo atual: (i) o \textit{full-batch GD} atualiza os pesos apenas uma vez a cada $56.000$ exemplos; (ii) a inicialização gaussiana fixa ($\sigma = 0{,}01$) não escala com o \textit{fan-in}, diferentemente da inicialização de He para ReLU; e (iii) o achatamento de imagens $28 \times 28$ em vetores de $784$ dimensões descarta a topologia espacial 2D, tornando o MLP sensível a translações, rotações e variações de espessura de traço em entradas desenhadas à mão.
